
\subsection*{Cônicas Rotacionadas}

\begin{frame}[label=conicas]{Cônicas Rotacionadas}
Os termos mistos da forma $xy$ de uma equação de cônica podem ser 
eliminados através de um processo de diagonalização, que equivale a uma 
rotação dos eixos.
\medskip

Uma expressão da forma
\[{\color{blue}a}x^2+{\color{blue}b}y^2+{\color{blue}c}xy\]
é chamada de \dt{forma quadrática} em $x$ e $y$ e pode ser representada 
usando matrizes da seguinte forma:
\[{\color{blue}a}x^2+{\color{blue}b}y^2+{\color{blue}c}xy=
\begin{bmatrix}
x & y
\end{bmatrix}
{\color{blue}\begin{bmatrix}
a & c/2\\ c/2 & b
\end{bmatrix}}
\begin{bmatrix}
x \\ y
\end{bmatrix}= X^t {\color{blue}A} X.
\]

\end{frame}


\begin{frame}[label=conicas]{Teorema dos Eixos Principais}
Como ${\color{blue}A}$ é uma matriz simétrica, pelo Teorema Espectral, 
existe uma matriz ortogonal que diagonaliza ${\color{blue}A}$, isto é, 
\[Q^t {\color{blue}A} Q={\color{red}D},\]
onde ${\color{red}D}$ é uma matriz diagonal formada por autovalores de 
${\color{blue}A}$.
Fazendo a \dt{mudança de variáveis} 
\[X=Q\bar{X}\ \text{ ou, equivalentemente, }\  \bar{X}=Q^tX,\]
temos
\begin{equation*}
X^t {\color{blue}A} X
=(Q\bar{X})^t {\color{blue}A}(Q\bar{X})
=\bar{X}^t(Q^t {\color{blue}A}Q) \bar{X}
=\bar{X}^t{\color{red}D}\bar{X}
\end{equation*}



\end{frame}

\begin{frame}[label=conicas]{Teorema dos Eixos Principais}
Assim, se ${\color{red}\lambda_1,\lambda_2}$ são os autovalores de ${\color{blue}A}$ e 
\[\bar{X}=
\begin{bmatrix}
\bar{x} \\ \bar{y},
\end{bmatrix}\]
então
\begin{align*}
{\color{blue}a}x^2+{\color{blue}b}y^2+{\color{blue}c}xy& =X^t {\color{blue}A} X\\
& =\bar{X}^t{\color{red}D}\bar{X}\\
& = {\color{red}\lambda_1}\bar{x}^2+{\color{red}\lambda_2}\bar{y}^2,
\end{align*}
uma forma quadrática sem termos mistos.

\end{frame}


\begin{frame}[label=conicas]
\begin{exe}
Identifique e faça um esboça das seguintes cônicas:
\begin{enumerate}
\item $5x^2+4xy+2y^2=6$.
\item $2x^2-4xy-y^2+8=0.$
\item $x^2+4\sqrt{2}xy-y^2+2\sqrt{6}x+2\sqrt{3}y=0$
\end{enumerate}
\end{exe}
\end{frame}


\begin{frame}[label=conicas]
\begin{casa}
Identifique e faça um esboço das seguintes cônicas

\begin{enumerate}
\item  $5x^2-4xy+8y^2-36=0$
\item $11 x^{2} + 24 x y + 4 x + 4 y^{2} + 3 y + \frac{15}{4} = 0$
\end{enumerate}
\end{casa}
\end{frame}

\begin{frame}[label=conicas]



\end{frame}